\documentclass[11pt]{scrlttr2}
\usepackage[ngerman]{babel}

% ---------- Absender ----------
\setkomavar{fromname}{Absender Name}
\setkomavar{fromaddress}{Straße 1\\12345 Stadt}
\setkomavar{fromemail}{absender@example.com}
\setkomavar{subject}{Betreff des Briefes}
\setkomavar{signature}{Absender Name}
\setkomavar{date}{\today}

\begin{document}

\begin{letter}{%
    Empfänger Name\\
    Empfänger Straße 2\\
    54321 Empfängerstadt%
}

\opening{Sehr geehrte Damen und Herren,}

hier steht der Brieftext. Absätze werden durch Leerzeilen getrennt und
KOMA-Script setzt Anschrift, Betreff und Datum automatisch an die richtige
Stelle.

Mit freundlichen Grüßen

\closing{Mit freundlichen Grüßen}

\ps{PS: Ein Postskriptum ist optional.}

\end{letter}

\end{document}
